\section{一个地址怎样到达唯一器件}

若把处理器的地址线同时接到启动存储和 RAM，而不给出选择规则，两个器件可能同时驱动返回
数据，处理器也无法判断结果来自哪里。机器需要一个位于请求发起者和目标器件之间的地址
互连。它读取地址的高位，与预先接好的范围比较，只允许匹配的目标接收本次请求。

\subsection{一次事务包含哪些信息}

处理器访问外部状态时发出一份请求。教学互连中的请求至少含有以下字段：

\begin{lstlisting}
request:
  address       32-bit physical address
  operation     READ or WRITE
  width         1, 2, or 4 bytes
  write_data    valid only for WRITE
  byte_enable   which byte lanes may change
\end{lstlisting}

目标完成后返回：

\begin{lstlisting}
response:
  status        OK, UNMAPPED, ALIGNMENT, or READ_ONLY
  read_data     valid only when a READ completes with OK
\end{lstlisting}

请求和响应可能由并行导线、片上网络分组或多周期握手承载。本书把从请求被接受到响应返回的
这段交互称为总线事务（bus transaction）。这里的“总线”描述软件可观察的传送约定，不
要求真实电路采用某个特定工业协议。

宽度字段不能省略。同一地址可用于读取一个字节或四个字节，目标必须知道哪些字节参与本次
操作。写请求还需要字节使能，避免更新相邻位置。读请求中的写数据没有意义，写请求返回的
读数据也没有意义；明确这些有效条件可以防止把任意电平误当成软件结果。

\subsection{地址互连怎样进行选择}

本样机采用一张固定物理地址表：

\begin{longtable}{|L{0.30\linewidth}|L{0.25\linewidth}|L{0.33\linewidth}|}
\hline
\rowcolor{BookTableHead}
物理地址范围 & 被选中的目标 & 当前用途 \\ \hline
\code{0x00000000--0x0000ffff} & 启动存储 & 启动程序与复位入口 \\ \hline
\code{0x00100000--0x001fffff} & RAM & 人工装入的程序、数据、栈与结果 \\ \hline
其余地址 & 无目标 & 返回 \code{UNMAPPED} \\ \hline
\end{longtable}

例如地址 \code{0x00102004} 落在 RAM 范围内。互连从 RAM 起始地址中减去
\code{0x00100000}，得到器件内部偏移 \code{0x00002004}，再把请求转发给 RAM。启动存储
在这次事务中保持不响应。若地址为 \code{0x20000000}，没有范围匹配，互连
直接返回 \code{UNMAPPED}，不能让处理器无限等待。

\begin{center}
\begin{tikzpicture}[node distance=8mm and 5mm]
\node[stage,minimum width=2.55cm] (req) {请求锁存器\\固定全部请求字段};
\node[stage,minimum width=2.7cm,right=of req] (cmp) {并行范围比较\\启动区、RAM 区};
\node[stage,minimum width=2.8cm,right=of cmp] (sel) {唯一选择与等待记录\\\code{selected=RAM}};
\node[stage,minimum width=3.0cm,right=of sel] (target) {目标请求端口\\只使能 RAM 通道};
\node[stage,minimum width=2.6cm,below=13mm of target] (tresp) {目标响应端口\\仅接收 RAM 响应};
\node[stage,minimum width=2.65cm,left=of tresp] (mux) {响应选择器\\核对等待目标};
\node[stage,minimum width=2.55cm,left=of mux] (out) {响应锁存器\\返回处理器};
\node[stage,fill=BookWarm,minimum width=2.7cm,left=of out] (err) {零匹配或多匹配\\本地错误响应};
\draw[flow] (req) -- (cmp);
\draw[flow] (cmp) -- (sel);
\draw[flow] (sel) -- (target);
\draw[flow] (target) -- (tresp);
\draw[flow] (tresp) -- (mux);
\draw[flow] (mux) -- (out);
\draw[->,>=Latex,thick,BookMuted] (cmp.south) |- (err.north);
\draw[->,>=Latex,thick,BookMuted] (err) -- (out);
\end{tikzpicture}
\end{center}
\diagramnote{地址 \code{0x00102004} 使 RAM 比较结果为 1。互连把这一位选择和请求字段一起
保持到响应返回，因此较晚到达的响应仍能被送回正确事务；零匹配或多匹配由互连直接报错。}

\subsection{地址互连的三层含义}

\textbf{硬件模型。} 互连保存或固化一组地址范围，接收请求字段，产生目标选择信号，再把
唯一目标的响应送回发起者。没有匹配或出现重叠匹配时，它产生错误响应。

\textbf{软件模型。} 程序看到统一的物理地址空间。相同的 \code{LOAD}/\code{STORE}
指令访问不同范围时会到达不同器件；程序必须遵守每个范围允许的宽度、方向和对齐规则。

\textbf{抽象。} 软件可以用一个物理地址唯一指出当前机器中的一个可寻址位置。这个能力
不包含虚拟内存、名称、访问权限或并发仲裁；它只回答“这次请求交给谁”。

图中的“等待记录”不能由组合比较线替代。请求被接受以后，处理器可能撤掉原地址，RAM
也可能经过若干周期才响应；互连仍要记住本次选择的是 RAM，并在响应返回前拒绝另一笔会
覆盖记录的请求。当前单发起者互连至少保存 \code{busy}、\code{selected\_target} 和
必要请求字段。响应提交后，这些字段才清除。

\subsection{一次 RAM 读取的完整时序}

以第一轮循环读取整数 \(7\) 为例。执行 \code{LOAD r5,[r0+0]} 时，
\code{r0=0x00102000}。处理器先形成地址，再等待目标完成：

\begin{enumerate}
\item 处理器计算有效地址 \code{0x00102000+0}；
\item 它发出 4 字节读请求，并在请求未完成时保留该条指令的上下文；
\item 互连选择 RAM，把地址换算为器件内部偏移 \code{0x00002000}；
\item RAM 读取连续四个字节 \code{07 00 00 00}；
\item RAM 返回 \code{OK} 和读数据，互连把响应送回处理器；
\item 处理器把组合后的 32 位值 \code{0x00000007} 写入 \code{r5}；
\item 只有写回发生后，PC 才按该指令的完成规则进入下一位置。
\end{enumerate}

若 RAM 需要多个时钟才能响应，处理器不能提前把 \code{r5} 改成未知值，也不能把这条
装入当作已经完成。教学处理器在等待期间暂停该指令的提交。更复杂处理器可以执行其他独立
工作，但最终仍必须使软件观察到与规定顺序一致的结果；本章不需要展开内部优化。

\subsection{时钟与复位怎样把随机上电状态收束到第一条取指}

存储和互连已经接好，处理器仍需要一个共同的状态更新节拍，并且需要在上电后得到确定的
起点。电源刚稳定时，各寄存器中的位不应被假定为零。复位电路必须主动把少量关键状态设置
为规定值，其中最重要的是 PC 和控制状态。

\subsection{时钟不是软件计时器}

时钟源周期性地产生边沿，时序电路在边沿处提交新状态。它使“先读取旧寄存器、再写入新
寄存器”的顺序有了物理边界，但普通软件此时没有可读的时钟计数器。知道电路有时钟，不等于
能够请求“十毫秒后通知我”。可编程计时器要在后续确有需求时作为新器件加入。

\textbf{硬件模型。} 时钟源输出周期边沿；复位电路在电源和时钟尚未稳定时保持复位有效，
稳定后按规定边沿释放。处理器在复位有效期间不提交普通指令状态。

\textbf{软件模型。} 复位释放后，PC 被设为 \code{0x00000100}，SP 和通用寄存器的值除
非另有规定，否则视为未定义。固件不能在初始化前读取这些寄存器并假定它们为零。

\textbf{抽象。} 启动代码可以依赖每次复位都从同一入口开始，并且已提交的指令状态在时钟
边界上遵守体系结构规则。它不能依赖真实频率恒定，也不能由边沿数推断日期或调度时刻。

时钟源内部也不是一个只写着“时钟”的方框。振荡单元先产生连续变化的原始信号；整形级把
缓慢或带噪的过渡收束为数字高低电平；可选分频状态决定输出边沿相隔多少个原始周期；输出
缓冲再把同一时钟送到处理器、互连的事务状态和 RAM 控制端。第一章只依赖这些部件看到同一
规定边沿，不依赖振荡器的材料、真实频率或相位调节方法。

\begin{center}
\begin{tikzpicture}[node distance=9mm and 7mm]
\node[stage,minimum width=2.7cm] (osc) {振荡单元\\产生原始周期信号};
\node[stage,minimum width=2.7cm,right=of osc] (shape) {电平整形\\形成明确高低电平};
\node[stage,minimum width=2.7cm,right=of shape] (divide) {分频状态\\按固定比率选取边沿};
\node[stage,minimum width=2.75cm,right=of divide] (buffer) {输出缓冲与扇出\\形成共同更新边沿};
\node[stage,minimum width=2.7cm,below=12mm of buffer] (receivers)
  {时序状态接收端\\处理器、互连、RAM 控制};
\draw[flow] (osc) -- (shape);
\draw[flow] (shape) -- (divide);
\draw[flow] (divide) -- (buffer);
\draw[flow] (buffer) -- (receivers);
\end{tikzpicture}
\end{center}
\diagramnote{分频状态和输出缓冲属于时钟路径本身；它们不构成软件可读计数器。处理器只能
依赖状态在规定边沿更新，尚不能请求某个未来时刻的通知。}

复位并不是一根从电源直接连到 PC 的线。电源监测器先给出“供电可用”，时钟稳定检测器再
确认已有连续边沿；复位保持状态随后等待规定边沿数，并通过同步释放级只在时钟边界改变
\code{reset\_active}。处理器中的复位选择器据此把 PC 和控制阶段写成规定起点，其他未
承诺清零的寄存器不经过这条初始化通路。

\Needspace{0.25\textheight}
\begin{center}
\begin{tikzpicture}[node distance=8mm and 4mm]
\node[stage,minimum width=2.3cm] (power) {电源监测器\\供电是否可用};
\node[stage,minimum width=2.35cm,right=of power] (clock) {时钟稳定检测\\是否已有连续边沿};
\node[stage,minimum width=2.35cm,right=of clock] (hold) {复位保持状态\\计满规定边沿};
\node[stage,minimum width=2.35cm,right=of hold] (sync) {同步释放级\\复位保持或释放};
\node[stage,minimum width=2.55cm,right=of sync] (init) {处理器复位选择\\PC、控制阶段取初值};
\draw[flow] (power) -- (clock);
\draw[flow] (clock) -- (hold);
\draw[flow] (hold) -- (sync);
\draw[flow] (sync) -- (init);
\end{tikzpicture}
\end{center}
\diagramnote{复位保持和同步释放把模拟上电过程收束为一个数字状态变化。处理器只对契约明确
列出的状态采用复位初值；RAM 和普通通用寄存器不会因此自动全部清零。}

\subsection{从电源稳定到第一条指令完成}

下面的时间线把“开机”拆成可观察事件。纵向虚线表示时钟边沿，横向各行表示不同部件：

\begin{center}
\begin{tikzpicture}[x=1.45cm,y=0.88cm]
\foreach \x/\lab in {1/{边沿 1},2/{边沿 2},3/{边沿 3},4/{边沿 4},5/{边沿 5}} {
  \draw[dashed,BookRule] (\x,0.25) -- (\x,4.75);
  \node[font=\scriptsize,BookMuted] at (\x,0) {\lab};
}
\node[anchor=east,font=\footnotesize] at (0.65,4.25) {电源};
\node[anchor=east,font=\footnotesize] at (0.65,3.25) {复位};
\node[anchor=east,font=\footnotesize] at (0.65,2.25) {PC};
\node[anchor=east,font=\footnotesize] at (0.65,1.25) {互连};
\fill[BookTint] (0.75,3.9) rectangle (1.2,4.6);
\fill[BookAccent!12] (1.2,3.9) rectangle (5.35,4.6);
\draw[BookRule] (0.75,3.9) rectangle (5.35,4.6);
\node[font=\scriptsize] at (0.98,4.25) {建立};
\node[font=\scriptsize] at (3.25,4.25) {稳定};
\fill[BookWarm] (0.75,2.9) rectangle (3,3.6);
\fill[BookTint] (3,2.9) rectangle (5.35,3.6);
\draw[BookRule] (0.75,2.9) rectangle (5.35,3.6);
\node[font=\scriptsize] at (1.85,3.25) {保持有效};
\node[font=\scriptsize] at (4.15,3.25) {已释放};
\fill[BookTint] (0.75,1.9) rectangle (3,2.6);
\fill[BookAccent!12] (3,1.9) rectangle (4,2.6);
\fill[BookTint] (4,1.9) rectangle (5.35,2.6);
\draw[BookRule] (0.75,1.9) rectangle (5.35,2.6);
\node[font=\scriptsize] at (1.85,2.25) {尚不可用于取指};
\node[font=\scriptsize] at (3.5,2.25) {\code{0x00000100}};
\node[font=\scriptsize] at (4.68,2.25) {后继 PC};
\fill[BookTint] (0.75,0.9) rectangle (3,1.6);
\fill[BookAccent!12] (3,0.9) rectangle (4.1,1.6);
\fill[BookTint] (4.1,0.9) rectangle (5.35,1.6);
\draw[BookRule] (0.75,0.9) rectangle (5.35,1.6);
\node[font=\scriptsize] at (1.85,1.25) {空闲};
\node[font=\scriptsize] at (3.55,1.25) {第一笔取指};
\node[font=\scriptsize] at (4.72,1.25) {下一笔事务};
\draw[very thick,BookAccent] (1.2,3.82) -- (1.2,4.68);
\draw[very thick,BookAccent] (3,0.82) -- (3,3.68);
\draw[very thick,BookAccent] (4,1.82) -- (4,2.68);
\end{tikzpicture}
\end{center}
\diagramnote{电源稳定并不立即开始取指。复位先保持若干边沿，释放后才把规定的 PC 值用于
第一笔事务；第一条指令完成后，PC 才进入后继值。}

第一条指令位于启动存储，而不是 RAM 中的任务入口。原因很直接：RAM 此时尚未装入任务，
输入参数和栈也未建立。复位 PC 若直接指向 \code{0x00100000}，处理器会把上电后的不确定
RAM 字节当作指令，机器行为无法预测。

\subsection{处理器怎样把指令字节变成状态变化}

处理器是当前唯一能够主动发起取指和数据事务的部件。它内部至少要保留 PC、SP、通用
寄存器、标志和当前指令的执行控制状态。这里不展开门级数据通路，只说明后续软件真正能
依赖的边界。

\textbf{硬件模型。} 处理器根据 PC 发出取指请求，译码返回的操作码与字段，读取旧寄存器，
执行算术或形成访存地址，并在指令完成时提交目标寄存器、内存效果和下一 PC。访存错误会
阻止普通写回，并把失败位置与原因留在处理器内部的停止记录中。

\textbf{软件模型。} 软件使用本章定义的寄存器和指令。除分支、调用、返回和显式跳转外，
PC 每完成一条固定长指令增加 4。调用约定规定哪些寄存器传参、SP 如何移动以及返回值放在
哪里，这些规则不是电路自动理解的类型系统。

\textbf{抽象。} 软件得到一条按指令集规则推进的执行流：每条完成的指令都有确定的前置
状态和后置状态。当前处理器不隔离任务，不保存多个执行现场，也不会因为另一个任务等待而
主动停止当前任务。

\subsection{第一条取指逐字段走读}

复位释放后，PC 为 \code{0x00000100}。处理器发出的请求为：

\begin{lstlisting}
address     = 0x00000100
operation   = READ
width       = 4
write_data  = ignored
byte_enable = 1111
\end{lstlisting}

互连发现地址属于启动存储。这四个字节就是启动程序的第一条普通指令；本章样机把它规定为
读取人工准备记录中提交标志的 \code{LOAD}。启动存储返回数据及 \code{OK} 后，处理器
完成译码并等待相应的数据读取；只有该装入指令成功，目标寄存器和顺序后继
\code{0x00000104} 才一同提交。这里不需要一条尚未解释的特殊“启动跳转”：复位 PC
本身就指向启动程序入口。

\begin{center}
\begin{tikzpicture}[node distance=11mm]
\node[stage,minimum width=3.1cm] (cpu1) {处理器\\用 PC 形成地址};
\node[stage,minimum width=3.1cm,right=of cpu1] (bus) {地址互连\\选择启动存储};
\node[stage,minimum width=3.1cm,right=of bus] (rom) {启动存储\\返回四个字节};
\node[stage,minimum width=3.1cm,right=of rom] (cpu2) {处理器\\译码并提交新 PC};
\draw[flow] (cpu1) -- node[above,font=\scriptsize]{读请求} (bus);
\draw[flow] (bus) -- node[above,font=\scriptsize]{选中的请求} (rom);
\draw[flow] (rom) -- node[above,font=\scriptsize]{数据+\code{OK}} (cpu2);
\node[font=\scriptsize,BookMuted,below=7mm of cpu2,align=center]
  {提交完成以后\\新 PC 用于下一笔取指};
\end{tikzpicture}
\end{center}
\diagramnote{取指不是处理器内部凭空取得指令。它是一笔带地址的读取事务；只有响应成功并
完成译码后，处理器才提交新的 PC。}

\subsection{顺序 PC、分支 PC 与返回 PC}

求和任务中出现三类下一 PC：

\begin{enumerate}
\item 普通指令完成后，下一 PC 为当前 PC 加 4；
\item \code{BNEZ} 根据 \code{r1} 的值，在循环入口和顺序后继之间选择；
\item \code{RET} 从 SP 指向的内存中读取返回地址，再把该值写入 PC。
\end{enumerate}

三者最终都只是决定 PC 的新值，但所依赖的状态不同。普通前进只依赖当前 PC；条件分支还
依赖寄存器或标志；返回还要成功完成一次栈内存读取。把它们都写成“跳到下一条”会隐藏
错误来源。

\subsection{把已经解释的连接重新合在一起}

章首先给出的整体图用于确定方位。现在，每个方框和箭头都已有明确职责，可以用同一幅结构
复核一次请求怎样往返。机器仍然只包含时钟与复位、处理器、地址互连、启动存储和 RAM。

\begin{center}
\begin{tikzpicture}[
  node distance=11mm and 18mm,
  hw/.style={stage,minimum width=3.4cm},
  sig/.style={->,>=Latex,draw=BookMuted,thick},
  request/.style={->,>=Latex,thick,draw=BookAccent},
  response/.style={->,>=Latex,thick,draw=BookMuted}
]
\node[hw] (cpu) {处理器\\PC、SP、寄存器};
\node[hw,right=of cpu,minimum width=3.7cm] (bus) {地址互连\\译码与响应返回};
\node[hw,above right=13mm and 21mm of bus] (rom) {启动存储\\固件};
\node[hw,below right=13mm and 21mm of bus] (ram) {RAM\\任务与运行状态};
\node[hw,above=14mm of cpu] (clock) {时钟与复位\\更新边界与起点};
\draw[sig] (clock) -- node[right,font=\footnotesize]{边沿/复位} (cpu);
\draw[request] ([yshift=3mm]cpu.east) -- node[above,font=\footnotesize]{请求} ([yshift=3mm]bus.west);
\draw[response] ([yshift=-3mm]bus.west) -- node[below,font=\footnotesize]{响应} ([yshift=-3mm]cpu.east);
\draw[request] ([yshift=4mm]bus.east) -| node[above,font=\scriptsize,pos=0.72]{选中的请求}
  ([yshift=3mm]rom.west);
\draw[response] ([yshift=-3mm]rom.west) |- node[below,font=\scriptsize,pos=0.28]{响应}
  ([yshift=1mm]bus.east);
\draw[request] ([yshift=-1mm]bus.east) -| node[above,font=\scriptsize,pos=0.72]{选中的请求}
  ([yshift=3mm]ram.west);
\draw[response] ([yshift=-3mm]ram.west) |- node[below,font=\scriptsize,pos=0.28]{响应}
  ([yshift=-4mm]bus.east);
\end{tikzpicture}
\end{center}
\diagramnote{此时机器已经能从固定启动程序取指并访问 RAM。地址互连只负责按地址送达，
不负责判断 RAM 中的字节是否构成一份完整程序。}

这幅图说明了当前能力的上限。复位能找到启动程序，启动程序能读写 RAM；但本章仍需由
操作人员在复位期间预置求和程序。机器运行后没有接收新程序的通道，也没有谁替操作人员
检查一串外部字节是否完整。
